% ===== ANALYSIS =====
% ===== ANALYSIS EXECUTION NOTES =====
% This section should answer "why main results happen", not add new primary claims.
% Keep the following mapping:
%   Fig-Diagonal (required):
%     PR^l(attn) vs mean diagonal weight, with linear fit and R^2.
%     file: figures/paper/fig_diagonal_vs_pr.pdf
%
%   Fig-Attn-Maps (supporting):
%     baseline / XSA / Gated-XSA map comparison for concentration pattern.
%     file: figures/paper/fig_attn_maps_triplet.pdf
%
%   Optional (appendix if space is tight):
%     "compute fraction W^l" plot and expanded layerwise diagnostics.
%
% Structure (revised 2026-04-27):
%   §5.1 Why is the parallel component structurally large? (theoretical)
%   §5.2 Quantifying cheap compute inside expensive operations
%   §5.3 Parallel component ↔ diagonal: empirical validation
%   §5.4 XSA's diagonal over-concentration: characterization

\section{Analysis}
\label{sec:analysis}

\subsection{Why Is the Parallel Component Structurally Large?}
\label{sec:why}

The prevalence of the parallel component is not a learned artifact—it
follows structurally from the pre-norm residual design.  At layer $l$, the
residual norm accumulates with depth:
\begin{equation}
  \|\hl{l}\| \approx \|\hl{0}\| + \sum_{k<l} \|\Dl{k}\|,
\end{equation}
while the block output $\Dl{l} = f^l(\LN(\hl{l}))$ is computed on a
\emph{normalized} input with $\|\LN(\hl{l})\| \approx \sqrt{d}$,
bounding $\|\Dl{l}\|$ independently of depth.  As $\|\hl{l}\|$ grows and
$\|\Dl{l}\|$ remains bounded, the residual update becomes a progressively
smaller perturbation of a large vector.  The $\LN$ operation removes the
magnitude of $\hl{l}$ before passing it to $f^l$, so the block's output
retains a direction correlated with the normalized input—which is aligned
with $\hl{l}$ itself.  Consequently, $\PR{l}$ tends to increase with depth
regardless of block type.

Formally, if we model $\Dl{l}$ as a random vector with fixed magnitude
$\epsilon$ and direction drawn uniformly from the unit sphere in $\mathbb{R}^d$,
conditioned on the $\LN$ normalization, the expected parallel ratio is:
\begin{equation}
  \mathbb{E}[\pr{l}{t}]
    = \mathbb{E}\!\left[\frac{|\Dl{l} \cdot \hat{\hl{l}}|}{\|\Dl{l}\|}\right]
    \approx \frac{1}{\sqrt{d}},
\end{equation}
\noindent
growing with depth as $\|\hl{l}\|$ increases and $f^l$'s output direction
correlates more strongly with $\hl{l}$ through the LN normalization path.
This toy-model argument is not intended as a strict theorem for trained
transformers, but as a structural explanation for why high $\PR{l}$ is
expected under pre-norm residual accumulation.

This argument applies equally to $\Attn$ and $\FFN$ blocks, explaining the
universality of the phenomenon observed in Figure~\ref{fig:profiles}.

\subsection{Cheap Compute Inside Expensive Operations}
\label{sec:waste}

To quantify the cost asymmetry, we define the \emph{parallel compute
fraction} at layer $l$ as
\begin{equation}
  W^l = \PR{l} \times \mathrm{FLOPs}(f^l),
\end{equation}
the portion of the block's FLOP budget that produces an update equivalent
to a single scalar per token.  Figure~\ref{fig:waste} shows $W^l$ across
depth for $\Attn$ and $\FFN$ blocks.

\begin{figure}[t]
  \centering
  % Replace with:
  % \includegraphics[width=0.9\linewidth]{figures/paper/fig_parallel_compute_fraction.pdf}
  \fbox{\rule{0pt}{4.5cm}\rule{0.85\linewidth}{0pt}}
  \caption{Parallel compute fraction $W^l / \mathrm{FLOPs}(f^l)$ vs.\ layer
           depth, for $\Attn$ (blue) and $\FFN$ (orange), averaged over [N]
           pretrained models.  Shaded region: std.\ across models.
           Both grow monotonically with depth; $\Attn$ is consistently higher.}
  \label{fig:waste}
\end{figure}

In deeper layers, the estimated parallel compute fraction becomes a substantial
portion of both Attn and FFN budgets, indicating that a nontrivial part of
expensive block computation is devoted to direction-preserving updates.  This
supports the practical motivation for parallel control: it targets a large
compute share while preserving most downstream behavior.

\subsection{Parallel Component and Attention Diagonal: Empirical Validation}
\label{sec:diagonal-validation}

% NOTE: this is the key mechanism validation.
% Minimal required stats in caption/text:
%   slope, intercept, R^2, number of points (model-layer pairs).
% Recommended robustness (appendix):
%   per-model R^2 distribution and confidence intervals.
Section~\ref{sec:value-level} predicts that the attention parallel ratio scales
linearly with the mean diagonal attention weight $\mathbb{E}[\Adiag{t}]$.
Figure~\ref{fig:diagonal} plots $\PR{l}$ (Attn-only) against
$\mathbb{E}_{x,t}[\Adiag{t}]$ across all layers and models.

\begin{figure}[t]
  \centering
  % Replace with:
  % \includegraphics[width=0.9\linewidth]{figures/paper/fig_diagonal_vs_pr.pdf}
  \fbox{\rule{0pt}{4.5cm}\rule{0.85\linewidth}{0pt}}
  \caption{Attention parallel ratio $\PR{l}$ vs.\ mean diagonal attention
           weight $\mathbb{E}_{x,t}[\Adiag{t}]$ per layer across [N] models.
           Each point is one (model, layer) pair.  Line: linear fit ($r^2 =$ [VALUE]).}
  \label{fig:diagonal}
\end{figure}

The strong linear relationship ($r^2 = $ [VALUE]) confirms our derivation:
layers where tokens attend heavily to themselves are precisely those that
produce the largest parallel residual updates.  This validates the causal
chain: $\Adiag{t}$ large $\Rightarrow$ $\Dup{l}{t}$ large $\Rightarrow$
$\PR{l}$ large.

\subsection{\xsa's Off-Diagonal Over-Concentration}
\label{sec:xsa-analysis}

% NOTE: keep this subsection short in main paper.
% What must be shown:
%   diagonal suppression alone is not enough; distribution can become peaky off-diagonal.
% Suggested quantitative companion (table or inline stats):
%   top-k mass or off-diagonal entropy change (baseline -> XSA -> Gated-XSA).
While \xsa{} forces $\Adiag{t} = 0$ and thereby reduces $\PR{l}$
(explaining its pretraining gains), the pre-softmax masking introduces a
secondary effect: the softmax renormalization concentrates the remaining
probability mass on whichever off-diagonal positions had the highest raw
scores.  In practice, this leads to [DESCRIBE OVER-CONCENTRATION — \eg,
a few neighboring positions capturing most of the redistributed weight,
or certain heads collapsing to near-deterministic patterns over a fixed
offset].

Figure~\ref{fig:attn-maps} shows representative attention maps for
baseline, \xsa{}, and \method{}.

\begin{figure}[t]
  \centering
  % Replace with:
  % \includegraphics[width=0.95\linewidth]{figures/paper/fig_attn_maps_triplet.pdf}
  \fbox{\rule{0pt}{5cm}\rule{0.9\linewidth}{0pt}}
  \caption{Attention maps averaged over heads at layer $[L]$ for baseline
           (left), \xsa{} (center), and \method{} (right).
           \xsa{} eliminates the diagonal but concentrates mass on
           [DESCRIPTION OF PATTERN].
           \method{} achieves a more uniform distribution while keeping
           $\Adiag{t}$ suppressed.}
  \label{fig:attn-maps}
\end{figure}

\method{} addresses this by [EXPLANATION OF FIX], yielding attention maps
with both low diagonal weight and broadly distributed off-diagonal mass—the
combination that produces the lowest pretraining loss
(Table~\ref{tab:proposed}).
